\section{From passages to absence of conjugate points}
\label{sec:global}

The passage estimates concern the intrinsic index form of the original
field on an entire visit to the open neck.  We first show that these
estimates combine along every complete geodesic.  We then give the
index-form argument that excludes conjugate points.

\begin{proposition}[Assembly of passage estimates]
\label{prop:global-assembly}
Let $(M,g)$ be a complete smooth Riemannian manifold and let $U\subset M$
be open.  Suppose that $g$ has nonpositive sectional curvature at every
point of $M\setminus U$.  Assume also that, for every complete unit-speed
geodesic $\gamma:\mathbb R\to M$, every connected component $J$ of
$\gamma^{-1}(U)$, and every smooth compactly supported field $V$ along
$\gamma$, the intrinsic index form satisfies
\[
 I_J(V)\geq 0.
\]
Then $I_{\mathbb R}(V)\geq0$ for every such $\gamma$ and $V$.
The same conclusion follows if the passage hypothesis is given in the
stronger form $I_J^\perp(V)\geq0$.
\end{proposition}

\begin{proof}
Put $T=\dot\gamma$ and write
\[
 i_V(s)=|D_sV|^2-
       \langle R(V,T)T,V\rangle.
\]
This is a continuous, compactly supported function, so it is absolutely
integrable.  The open subset $\gamma^{-1}(U)$ of $\mathbb R$ is a
countable disjoint union of open intervals $J_j$: each interval contains
a distinct rational number.  Its complement $E$ is closed and hence
measurable.  Countable additivity gives
\begin{equation}
 I_{\mathbb R}(V)=\sum_j I_{J_j}(V)+\int_E i_V(s)\,ds,
 \qquad
 \sum_j|I_{J_j}(V)|\leq\int_{\mathbb R}|i_V(s)|\,ds<\infty.
 \label{eq:countable-index-assembly}
\end{equation}
Every summand is nonnegative by hypothesis.  On $E$, nonpositive
sectional curvature implies $i_V\geq0$: the curvature term only depends
on the component of $V$ perpendicular to $T$.  This proves the claim.
Finally,
\[
 I_J(V)=I_J^\perp(V)
       +\int_J\left|\frac{d}{ds}\langle V,T\rangle\right|^2ds
\]
proves the last assertion.
\end{proof}

This argument allows infinitely many passages and accumulating passage
endpoints.  It also includes an interval equal to $\mathbb R$ and
intervals with one infinite endpoint.  At a finite endpoint, the
restriction of $V$ extends smoothly with arbitrary boundary data; at an
infinite endpoint it vanishes beyond a finite time.  These are precisely
the fields allowed in the passage estimates.  Any tangential changes
used to prove those estimates are made separately on each passage.
They need not fit together across endpoints: their normal index
integrands equal that of the original field, and
\eqref{eq:countable-index-assembly} sums only the original intrinsic
integrals.  In particular, no bounds uniform over the auxiliary gauge
fields are required.  The integral over $E$ includes all times on the
interfaces, even if that set has positive measure.

\begin{lemma}[The index-form criterion]
\label{lem:index-criterion}
Let $(M,g)$ be a complete smooth Riemannian manifold.  If
$I_{\mathbb R}(V)\geq0$ for every smooth compactly supported field along
every complete unit-speed geodesic, then $g$ has no conjugate points.
\end{lemma}

\begin{proof}
Fix a complete unit-speed geodesic $\gamma$, set $T=\dot\gamma$, and use
the symmetric bilinear index form
\[
 B(V,W)=\int_{\mathbb R}
   \bigl(\langle D_sV,D_sW\rangle
        -\langle R(V,T)T,W\rangle\bigr)\,ds.
\]
Thus $I(V)=B(V,V)$.  We first extend the assumed nonnegativity to
compactly supported $H^1$ fields.  On a bounded interval containing the
support in its interior, choose a parallel orthonormal frame along
$\gamma$.  The coefficients of such a field belong to $H^1$, and their
zero extensions can be approximated in $H^1$ by smooth compactly
supported functions, using convolution with a smooth mollifier.
The frame identifies $D_s$ with ordinary differentiation.  The
curvature operator $W\mapsto R(W,T)T$ is bounded on this fixed compact
interval, so $B$ is continuous in the $H^1$ norm.  Consequently
$I(V)\geq0$ for every compactly supported $H^1$ field.

Suppose now that a nonzero Jacobi field $J$ along $\gamma$ vanishes at
two times $a<b$.  The Jacobi equation is
\[
 D_s^2J+R(J,T)T=0.
\]
Let $\widehat J$ agree with $J$ on $[a,b]$ and vanish elsewhere.  Because
$J(a)=J(b)=0$, this extension is continuous and belongs to $H^1$; its
weak derivative equals $D_sJ$ on $(a,b)$ and zero outside.
Integration by parts and the Jacobi equation yield
\[
 I(\widehat J)
 =\bigl[\langle D_sJ,J\rangle\bigr]_{a}^{b}=0.
\]
For every smooth compactly supported field $W$ and every real $t$,
nonnegativity on $H^1$ fields gives
\[
 0\leq I(\widehat J+tW)
       =2tB(\widehat J,W)+t^2I(W).
\]
Considering both signs of sufficiently small $t$ shows that
$B(\widehat J,W)=0$.  On the other hand, a second integration by parts
gives
\begin{equation}
 B(\widehat J,W)
   =\langle D_sJ(b),W(b)\rangle
       -\langle D_sJ(a),W(a)\rangle.
 \label{eq:jacobi-derivative-jump}
\end{equation}
The vector $D_sJ(a)$ is nonzero: otherwise the Jacobi equation with
initial data $J(a)=D_sJ(a)=0$ would imply $J\equiv0$ by uniqueness.
Choose a smooth field $W$, supported in a small neighborhood of $a$
disjoint from $b$, with $W(a)=D_sJ(a)$.  For example, multiply the
parallel extension of this vector by a smooth bump equal to one at
$a$.  Equation~\eqref{eq:jacobi-derivative-jump} then gives
$B(\widehat J,W)=-|D_sJ(a)|^2<0$, a contradiction.

This argument includes arbitrary Jacobi fields.  In particular, the
tangential coefficient $\langle J,T\rangle$ has second derivative
zero, so two zeros force that coefficient to vanish identically.
Every nonconstant affinely parametrized geodesic reduces to unit speed
by a linear change of parameter.  Along a constant geodesic the Jacobi
equation is $J''=0$, which has no nonzero solution with two zeros.
\end{proof}

\begin{proof}[Proof of Theorem~\ref{thm:main}]
We apply these results to the constructed metric and the open neck
$U=(0,1)\times T^2$.  The manifold is closed, so its smooth metric is
complete.  The product cores and the matching end collars have
nonpositive sectional curvature, including their boundary tori.
Proposition~\ref{prop:passage-index} supplies the passage hypothesis for
every complete unit-speed geodesic, with the one choice of the small
construction parameter already made there.  Proposition~
\ref{prop:global-assembly} and Lemma~\ref{lem:index-criterion} therefore
show that the metric has no conjugate points.

Proposition~\ref{prop:no-npc} excludes every smooth nonpositively curved
metric on this same manifold. The construction in
Section~\ref{sec:topology} makes it closed, connected, and orientable,
and Section~\ref{sec:cometric} gives the required $C^\infty$ metric.
This proves both assertions of the theorem.
\end{proof}
